\documentclass[11pt,a4paper]{article}

% Basic formatting and mathematics.
\usepackage[T1]{fontenc}
\usepackage[utf8]{inputenc}
\usepackage{lmodern,microtype}
\usepackage[margin=1in]{geometry}
\usepackage{amsmath,amssymb,amsthm,mathtools}
\usepackage{graphicx,booktabs,tabularx,enumitem}
\usepackage{longtable,pdfpages}
\usepackage[colorlinks=true,allcolors=blue]{hyperref}

% Shortcuts you can change or extend.
\newcommand{\R}{\mathbb{R}}
\newcommand{\A}{\mathcal{A}}
\newcommand{\K}{\mathcal{K}}
\newcommand{\ind}{\mathbf{1}}
\newcommand{\std}{\mathrm{std}}
\newcommand{\stu}{\mathrm{stu}}
\DeclareMathOperator{\cov}{cov}
\DeclareMathOperator{\diag}{diag}
\DeclareMathOperator*{\argmin}{arg\,min}
\DeclarePairedDelimiter{\norm}{\lVert}{\rVert}

\newtheorem{proposition}{Proposition}[section]
\theoremstyle{definition}
\newtheorem{definition}[proposition]{Definition}
\theoremstyle{remark}
\newtheorem*{remark}{Remark}

\title{MultiQT notes}
\author{} % Add your name here if you wish.
\date{\today}

\begin{document}
\maketitle

\section{Notation}
\label{sec:notation}
\begingroup\fontsize{9.5}{11.2}\selectfont
Unless stated otherwise, fix one disease, location and forecast horizon.

\paragraph{Forecasts, loss and projection.}
Indices $s,t$ order reference weeks (\texttt{reference\_date}, Saturday);
$j$ can denote a historical step in a re-run. Outcome $y_t$ belongs to
the target week of $q_t$, identified by \texttt{target\_end\_date} in
the forecasts and \texttt{target\_date} in prepared cases.
For levels $\A=\{\alpha_1<\cdots<\alpha_m\}\subset(0,1)$, let
$b_t=(b_t^\alpha)_{\alpha\in\A}$ and $q_t=(q_t^\alpha)_{\alpha\in\A}$
be the base and recalibrated forecasts. Values are in outcome units
(hospital counts) unless a transformation is shown.
The methodological recursions forecast, observe $y_t$, then update;
the coding section replaces this timing by delayed or revised feedback.
With $x_+=\max(x,0)$ and indicator $\ind\{\cdot\}$, quantile loss is
\begin{equation}
 \rho_\alpha(z,r)=\alpha(r-z)_++(1-\alpha)(z-r)_+,
 \label{eq:quantile-loss}
\end{equation}
with prediction first and observation second. Assume $b_t\in\K$, where
\[
 \K=\{q\in\R^m:q^{\alpha_1}\le\cdots\le q^{\alpha_m}\},\qquad
 \Pi_C(x)=\argmin_{z\in C}\norm{z-x}^2.
\]
Here $C$ is nonempty, closed and convex, and $\norm{\cdot}$ is Euclidean.
For positive-definite $W$,
$\Pi_\K^W(v)=\argmin_{q\in\K}(q-v)^\top W(q-v)$.
$Q_p$ is an empirical $p$-quantile, $I$ the identity matrix, and
$\operatorname{clip}_{[a,b]}(x)=\min\{b,\max\{a,x\}\}$.

\paragraph{Scores and states.}
$S_t(y)$ is a score, $r_t=S_t(y_t)$ its observed value and $c_t$ its
threshold. Target coverage is $\alpha$, target miscoverage is
$\delta=1-\alpha$, and $\cov_t^\alpha=\ind\{y_t\le q_t^\alpha\}$.
Superscripts $\std,\stu$ distinguish standard and studentized methods.
A tilde marks hidden accumulated state; the untilded state is used
after projection. Offsets $\tilde\theta_t,\theta_t$ are in outcome
units in the theory and in the selected tracker units in the code;
$\tilde u_t,u_t$ are in studentized units. Positive finite scales form
$A_t=\diag(\sigma_t^{\alpha_1},\ldots,\sigma_t^{\alpha_m})$;
$\sigma_t$ is common to levels and $\sigma$ is constant.
The step is $\eta$ or $\eta_t$, and $\kappa$ is a rate multiplier
(\texttt{multiplier}, labelled $m$ in script 02, distinct from $m=|\A|$).
$g_\alpha$ is the selected subgradient in~\eqref{eq:score-gradient}.

\paragraph{Reports and feedback.}
$\iota_t$ is the information cutoff: normally the Wednesday of the
reference week, or the actual generation date where known.
$y_s^{[v]}$ is the latest report for the target week of $q_s$ dated on
or before cutoff date $v$ (archive field \texttt{version}). A valid
report is finite and nonnegative; $\nu_s$ is its first valid publication
date and $d$ the additional delay in calendar weeks.
$m_s^{(d)},\tau_d(s),\mathcal R_t^{(d)}$ are the maturity date,
consuming-round index and rows consumed before forecasting at $t$;
see~\eqref{eq:feedback-maturity}--\eqref{eq:feedback-release-set}.
$I_{20}^{(d)}$ contains the first 20 consumed rows in release order, with
\begin{equation}
 \bar y_s^{(d)}=y_s^{[\iota_{\tau_d(s)}]},\qquad \tau_d(s)\text{ exists}.
 \label{eq:initial-feedback-values}
\end{equation}
For forecasts and states, $[t]$ instead identifies the re-run at
reference-week index $t$: $q_s^{[t]}$ is its reconstructed forecast for
row $s$. Feedback follows the cutoff convention, with the retrospective
initialization qualification in Section~\ref{sec:tuning-initialization}.
$\mathcal F_t^{(d)}=\bigcup_{j\le t}\mathcal R_j^{(d)}$ contains all
consumed rows; $\mathcal F_{t,20}^{(d)}$ and $\mathcal H_t^{(d)}$
contain the last 20 and five, respectively, or fewer if unavailable.

\paragraph{Implementation.}
The tracker transform is $\psi(x)=x$ or $\log(1+x)$, with
$\beta_t^\alpha=\psi(b_t^\alpha)$ and centre $a_t^\alpha$ equal to
the full base or its median. $\Phi$ is the standard normal distribution
function. Base residuals are
$e_s^\alpha=\psi(\bar y_s^{(d)})-\beta_s^\alpha$; raw scale estimates
have a hat and bounded scales do not. The earlier paired-width method
uses $a_\alpha=\min(\alpha,1-\alpha)$, with the nearest lower grid
level at $.5$.
The count divisor $c_{\mathrm{scale}}$ is population/$100{,}000$ by
default; $c_{\mathrm{pop}}$ always uses that divisor.
$S_0$ (initial error scale) and $L_0$ (expert loss normalization), with
floors $\varepsilon_S,\varepsilon_L$, are defined in
\eqref{eq:initial-error-scale} and~\eqref{eq:expert-reference-loss}.
$S_0$ is distinct from the score function $S_t$.
Square-root benchmarks use transform $\chi$, projected forecast
$\xi_t$, fast/slow states $\tilde\theta_t^F,\tilde\theta_t^S$ and
released coverage-error sum $\epsilon_t^\alpha$.
The leaky expert has common log-scale shift $v_t$; $\Delta_t$ is elapsed
calendar weeks since the previous reference date, at least one.

\paragraph{Experts and weighting.}
$f_{k,t}^\alpha$ is the forecast of expert $k$. Update index $n$ counts
processed feedback, with forecast row $s(n)$ and report $Y_n$.
For $K$ experts, $w_{k,n}^\alpha$ and $\lambda_n^\alpha$ are the
pre-update weight and rate; $\bar w_{k,n+1}^\alpha$ is the weight after
the exponential step and before sharing. Common weights omit $\alpha$.
$\hat q$ is the mixture before crossing repair; time indices may be
suppressed for one update. Losses $L_{k,n}^\alpha$ are unnormalized,
$\ell_{k,n}^\alpha$ normalized; omitting $\alpha$ denotes their mean
over levels. Weighted variances are $M_n^\alpha$ or $M_n$, lookback is
$H$, starting rate is $\lambda_{\mathrm{start}}$, and sharing fraction
is $\gamma$. Superscripts $\mathrm{raw},\mathrm{norm}$ distinguish
calculations without and with $L_0$ normalization.
$K_{\mathrm{recent}}$ is the number of recent weekly starts in the
equal mixture, and $\mathcal B_t$ their set.
$S_{k,t}^\alpha(y)$ is the residual score of bank member $k$.

For the target signal, $Z(t,h;v)$ is $\log(1+\text{base median})$ at
reference date $t$ and horizon $h$, if finite, nonnegative, available
by $v$ and not filled; otherwise it is missing. Write
$Z_{t,h}=Z(t,h;\iota_t)$. Signal $D_{t,h}$ gives context
$c_{t,h}^{\mathrm{ctx}}\in\{0,1\}$ and gate $G_{t,h}^T\in[0,1]$,
the weight of the original in \texttt{score\_gated}. Its threshold
$\delta_G$ and width $w_T$ are distinct from miscoverage $\delta$.

\paragraph{Diagnostics and evaluation.}
For the report screen only, $u$ indexes target weeks,
$z_u^{[v]}=\log(1+y_u^{[v]})$, and
$d_u^{[v]}=z_u^{[v]}-z_{u-1}^{[v]}$; historical median, scaled MAD
and predicted log count are $\mu_u^{[v]},s_u^{[v]},\widehat z_u^{[v]}$.
For evaluation, $y_t^\star$ is the latest archived count, $T_t$ the
target date, $p(t)$ the nearest retained peak date (earlier at a tie),
and $u_t^{\mathrm{peak}}=(T_t-p(t))/7$.
$E$ contains evaluation rows with finite truth; $W$ is a selected
window and $N=|W|$. $\mathcal L_{j,c}$ and $\mathrm{CE}_{j,c}$ are
mean quantile loss and calibration error for method $j$ in case $c$;
$j=0$ is the original. A case fixes location, horizon and feedback
design within a disease; $\mathcal C$ is a set of $J$ cases.
Tuning candidate $\vartheta$ has mean loss
$\overline{\mathcal L}_j(\vartheta)$ over tuning cases, with
$\overline{\mathcal L}_0$ for the original.
For peak detection, $x_i$ are retained latest counts in date order,
$z_i$ their smoothed values, $P_i$ prominence, and $l_i,r_i$ the
nearest strictly higher smoothed positions (or series ends).
$X$ is the largest retained count on a forecast target date.
\endgroup

\clearpage
\section{Methodological notes}
\subsection{Studentized score}
Standard MultiQT tracks additive residual corrections. Studentization
tracks scaled residuals, so changing the scale also changes the
accumulated correction in outcome units. A falling scale can shrink a
positive correction after a peak without erasing the hidden state.
The quantile coverage target is unchanged; this construction alone
does not establish calibration guarantees.

\subsubsection{Standard MultiQT and score tracking}
Starting with zero hidden state, MultiQT~\cite{multiqt} projects the
forecast but updates the unprojected state:
\begin{align}
 \theta_t&=\Pi_{\K-b_t}(\tilde\theta_t),
 &q_t^\std&=b_t+\theta_t,\label{eq:standard-projection}\\
 \cov_t^{\std,\alpha}&=\ind\{y_t\le q_t^{\std,\alpha}\},
 &\tilde\theta_{t+1}^\alpha
 &=\tilde\theta_t^\alpha-\eta(\cov_t^{\std,\alpha}-\alpha).
 \label{eq:standard-update}
\end{align}
The scalar score tracker of Angelopoulos et al.~\cite{pid} uses
\begin{equation}
 \mathcal P_t=\{y:S_t(y)\le c_t\},\qquad
 \mathrm{err}_t=\ind\{r_t>c_t\},
 \label{eq:score-set}
\end{equation}
and updates
\begin{equation}
 c_{t+1}=c_t+\eta(\mathrm{err}_t-\delta)
 =c_t-\eta\bigl(\ind\{r_t\le c_t\}-\alpha\bigr).
 \label{eq:pid-update}
\end{equation}
This is a quantile-loss subgradient step. In the threshold argument,
\begin{equation}
 \partial_z\rho_\alpha(z,r)=
 \begin{cases}
 \{-\alpha\},&z<r,\\
 [-\alpha,1-\alpha],&z=r,\\
 \{1-\alpha\},&z>r,
 \end{cases}
 \label{eq:subgradient}
\end{equation}
and the selected subgradient, including ties, is
\begin{equation}
 g_\alpha(z;r)=\ind\{r\le z\}-\alpha
 \ \in\ \partial_z\rho_\alpha(z,r).
 \label{eq:score-gradient}
\end{equation}
Loss depends on distance from the threshold; the update depends only
on which side contains the observation.

\subsubsection{Residual coordinates}
The signed residual score recovers standard MultiQT:
\begin{equation}
 S_t^{\std,\alpha}(y)=y-b_t^\alpha,\qquad
 r_t^{\std,\alpha}=y_t-b_t^\alpha.
 \label{eq:standard-score}
\end{equation}
Its threshold $\theta_t^\alpha$ defines the forecast endpoint and coverage:
\begin{align}
 \{y:S_t^{\std,\alpha}(y)\le\theta_t^\alpha\}
 &=(-\infty,b_t^\alpha+\theta_t^\alpha]
 =(-\infty,q_t^{\std,\alpha}],\label{eq:residual-inversion}\\
 \cov_t^{\std,\alpha}
 &=\ind\{r_t^{\std,\alpha}\le\theta_t^\alpha\}.
 \label{eq:standard-coverage}
\end{align}
Thus
\begin{equation}
 \tilde\theta_{t+1}^\alpha=\tilde\theta_t^\alpha
 -\eta\bigl(\ind\{r_t^{\std,\alpha}\le\theta_t^\alpha\}-\alpha\bigr),
 \label{eq:residual-update}
\end{equation}
with the same loss in residual coordinates:
\begin{equation}
 \ell_t^\std(\theta)
 =\sum_{\alpha\in\A}\rho_\alpha(b_t^\alpha+\theta^\alpha,y_t)
 =\sum_{\alpha\in\A}\rho_\alpha(\theta^\alpha,r_t^{\std,\alpha}).
 \label{eq:standard-loss}
\end{equation}
For one level there is no ordering constraint, so this is scalar
quantile tracking. For several levels, evaluate the subgradient at the
projected state and apply it to the hidden state.

\subsubsection{Studentized residuals and their objective}
Choose positive scales from information available at the forecast
cutoff and hold them fixed when differentiating:
\begin{equation}
 S_t^{\stu,\alpha}(y)=\frac{y-b_t^\alpha}{\sigma_t^\alpha},\qquad
 r_t^{\stu,\alpha}=\frac{y_t-b_t^\alpha}{\sigma_t^\alpha}.
 \label{eq:student-score}
\end{equation}
The score is increasing, so inversion preserves the coverage event:
\begin{equation}
 q_t^{\stu,\alpha}=b_t^\alpha+\sigma_t^\alpha u_t^\alpha,\qquad
 \cov_t^{\stu,\alpha}=\ind\{r_t^{\stu,\alpha}\le u_t^\alpha\}.
 \label{eq:student-inversion}
\end{equation}
Quantile loss of the standardized scores is a weighted outcome-unit loss:
\begin{equation}
 \ell_t^\stu(u)=\sum_{\alpha\in\A}\rho_\alpha(u^\alpha,r_t^{\stu,\alpha})
 =\sum_{\alpha\in\A}\frac{
 \rho_\alpha(b_t^\alpha+\sigma_t^\alpha u^\alpha,y_t)}{\sigma_t^\alpha}.
 \label{eq:student-loss}
\end{equation}
Its selected subgradient is
\begin{equation}
 g_t^{\stu,\alpha}(u_t)
 =\ind\{r_t^{\stu,\alpha}\le u_t^\alpha\}-\alpha
 =\cov_t^{\stu,\alpha}-\alpha.
 \label{eq:student-gradient}
\end{equation}
Indeed, away from ties the inverse scale in the loss cancels the scale
from differentiation of the forecast:
\begin{equation}
 \frac{\partial}{\partial u^\alpha}
 \frac{\rho_\alpha(b_t^\alpha+\sigma_t^\alpha u^\alpha,y_t)}{\sigma_t^\alpha}
 =\ind\{r_t^{\stu,\alpha}\le u^\alpha\}-\alpha.
 \label{eq:scale-cancellation}
\end{equation}
At ties use~\eqref{eq:score-gradient}. Starting from $\tilde u_1=0$,
the update is
\begin{equation}
 \tilde u_{t+1}^\alpha=\tilde u_t^\alpha
 -\eta g_t^{\stu,\alpha}(u_t).
 \label{eq:student-update}
\end{equation}
The projected threshold determines coverage; only the hidden threshold
accumulates increments. This is not gradient descent on the unweighted
outcome-unit loss in $u$: that gradient has an additional factor
$\sigma_t^\alpha$. The coding section also tests median-centred forecasts
and nonzero initial states.

\subsubsection{Ordering and projection}
Ordering always concerns the returned forecasts. For standard residuals,
\begin{align}
 C_t^\std&=\K-b_t,\qquad b_t+\theta\in\K
 \ \Longleftrightarrow\ \theta\in C_t^\std,
 \label{eq:standard-constraint}\\
 \theta^{\alpha_i}-\theta^{\alpha_{i+1}}
 &\le b_t^{\alpha_{i+1}}-b_t^{\alpha_i},\qquad i=1,\ldots,m-1.
 \label{eq:standard-adjacent}
\end{align}
For studentized residuals,
\begin{equation}
 C_t^\stu=\{u:b_t+A_tu\in\K\}=A_t^{-1}(\K-b_t).
 \label{eq:student-constraint}
\end{equation}
This is a nonempty closed convex polyhedron containing zero, generally
not a cone. Unequal scales do not preserve $\K$, so it cannot generally
be written as $\K-A_t^{-1}b_t$. The unique Euclidean projection gives
\begin{equation}
 u_t=\Pi_{C_t^\stu}(\tilde u_t),\qquad q_t^\stu=b_t+A_tu_t.
 \label{eq:student-projection}
\end{equation}
Changing variables to $q=b_t+A_tu$ yields weighted isotonic regression:
\begin{equation}
 q_t^\stu=\argmin_{q\in\K}\sum_{\alpha\in\A}
 \frac{(q^\alpha-b_t^\alpha-\sigma_t^\alpha\tilde u_t^\alpha)^2}
 {(\sigma_t^\alpha)^2},
 \label{eq:weighted-projection}
\end{equation}
or equivalently,
\begin{equation}
 b_t+A_t\Pi_{A_t^{-1}(\K-b_t)}(\tilde u_t)
 =\Pi_\K^{A_t^{-2}}(b_t+A_t\tilde u_t).
 \label{eq:projection-identity}
\end{equation}
Weights $1/(\sigma_t^\alpha)^2$ follow from Euclidean projection in
score coordinates; ordinary residual coordinates give equal weights.

\subsubsection{Common and constant scales}
For $A_t=\sigma_t I$, positive scalar multiplication preserves $\K$:
\begin{equation}
 C_t^\stu=\K-b_t/\sigma_t,\qquad
 q_t^\stu=\Pi_\K(b_t+\sigma_t\tilde u_t).
 \label{eq:common-scale}
\end{equation}
The projection is ordinary isotonic regression. This is standard
MultiQT on $b_t/\sigma_t$ and $y_t/\sigma_t$, with the returned forecast
multiplied by $\sigma_t$. For constant $\sigma$, it is also equivalent
to standard MultiQT on the original inputs after matching initialization
and converting state and rate:
\begin{equation}
 \tilde\theta_t=\sigma\tilde u_t,\qquad
 \eta_\std=\sigma\eta_\stu.
 \label{eq:constant-common-scale}
\end{equation}
For varying $\sigma_t$, changing the scale rescales the accumulated
state; changing only the current rate does not reproduce that effect.

\subsection{Questions and ideas}
Which property of the base forecast should another score recalibrate?
Increasing, invertible scores preserve the quantile coverage event;
changing the target to PIT or highest-density-region calibration would
require a separate formulation.

\clearpage

\section{Coding Updates}
\subsection{Data}
\label{sec:data-preparation}
Script 00 downloads and generates base forecasts; script 01 prepares
cases from the CMU--TimeSeries hub quantiles and dated NHSN reports.
Cases retain generation cutoffs, revisions, feedback schedules for
delays 0--4, populations and filled-forecast flags. New data, tuned
settings and scores carry \texttt{latest-at-use-v2}; readers reject
missing or mismatched protocol tags.

\paragraph{Data designs.}
The active \texttt{hub\_first\_months} design uses both diseases from
23 November 2024. Target dates determine periods: tuning through
31 May 2025, then evaluation within the October--September seasons
2024/25 and 2025/26. Plotted weeks outside these seasons are not scored.
``Latest'' means the latest archived report, through 30 October 2026;
it does not imply that further revisions are impossible.

The \texttt{seasons} alternative starts on 14 October 2023, tunes on
2023/24 and evaluates on 2024/25 and 2025/26. Own COVID forecasts have
been generated with unchanged, pinned Delphi
\texttt{windowed\_\allowbreak seasonal} code and packages, using solver
\texttt{fn}. Flu still uses hub forecasts because own earlier forecasts
require non-public Delphi auxiliary data. This extended design has not
supplied the results below.

Report versions begin with a retrospective bulk load on 19 November 2024.
The earlier experiment therefore uses latest counts restricted to dates
then observable. Feedback through target week 27 April 2024 also uses
latest counts, with a synthetic Wednesday release after the target week
and flag \texttt{synthetic\_final}; historical revisions are unavailable.
Own horizon $-1$ forecasts have zero width because the target is already
a predictor. No forecasts span the May--November 2024 reporting gap;
stateful methods retain their state across it.

\paragraph{Cleaning and filling.}
\label{sec:data-issues}
Forecast cleaning removes submissions, without correcting reported counts.
\begin{enumerate}[leftmargin=*,itemsep=2pt]
\item \textbf{Missing forecasts.} Fill gaps of at most six rounds from
the most recent earlier forecast for the same target: horizon $h$ at
reference date $r$ uses $h+k$ at $r-7k$ days, with the smallest available
$k>0$. Horizon 3 has no source. Filled rows are flagged and scored;
unfillable rounds remain absent.
\item \textbf{Broken submissions.} Remove the known Delphi bugs:
COVID 23 November 2024; COVID $h=-1$ on 25 April 2026; US flu 14 and
21 December 2024 and 4 January 2025. Also remove all-zero submissions
across every location and quantile of a disease--round--horizon
(flu $h=-1$, 22 November 2025), then apply the same filling rule.
There is no general magnitude-based forecast filter.
\item \textbf{Missing reports.} If unpublished or currently invalid,
stateful feedback waits and re-runs omit it; neither substitutes an
older valid report. The active cohort has no stored missing values,
although a report can be unavailable at an earlier cutoff. Targets
4 May--2 November 2024 and forecast rows without finalized truth are
excluded from learning and scoring.
\item \textbf{Report errors.} The main workflow rejects nonfinite or
negative reports only. Finite, nonnegative errors can pass: Pennsylvania
COVID for 16 November 2024 changed from 6,124 to 1,005 in successive
releases. The screen below and withholding experiment are confined to
\texttt{agent/}; neither is adopted in \texttt{R/}.
\end{enumerate}

\begin{center}\small
\begin{tabular}{@{}lll@{}}
\toprule
Disease & Missing reference weeks & Unfilled target weeks, $h=2$\\
\midrule
COVID & 25 January--1 February 2025 & 15 February 2025\\
Both & 25 October--15 November 2025 & 15--29 November 2025\\
Flu & 25 January 2025; 18 April 2026 & None\\
Flu & 7 June--11 October 2025 & 21 June--25 October 2025\\
Flu & 6 June--3 October 2026 & 20 June--17 October 2026\\
\bottomrule
\end{tabular}
\end{center}
The COVID $h=2$ weight plots (page~\pageref{plot:weights}) show the first
two gaps; only the first missing round can use an earlier $h=3$ forecast.
US flu also lacks target 4 January 2025 at $h=2$ after successive broken
submissions. Lighter bars denote filled rows. Lines break at gaps longer
than 28 days and can bridge shorter missing intervals.
Updates occur only at retained rounds: pending feedback waits across
unfilled rounds, without resetting the state.

\paragraph{Locations.}
Every disease--location--horizon case learns independently. Detailed
plots show the US, California, Texas, New York and Pennsylvania; this
is the selected list, not the four largest states by population.
Scores cover all 53 locations. Tuning and pooled summaries use the
52 non-US locations; the US is evaluated separately with their selected
parameters. Forecasts are never averaged across locations.

\paragraph{Experimental report screen.}
\label{sec:report-screen}
The experiment is in \path{agent/scripts/lib/report_screen.R}, with
corrected results in \path{agent/outputs/report_screen_notes_audit/README.md}.
At cutoff $v$, select the latest version before checking validity.
For target week $u$, use up to eight preceding consecutive-week changes,
without bridging missing weeks. Given at least six valid changes and
a valid previous target week,
\begin{align}
 \mu_u^{[v]}&=\operatorname{median}_j d_j^{[v]},&
 s_u^{[v]}&=1.4826\operatorname{median}_j|d_j^{[v]}-\mu_u^{[v]}|,\notag\\
 \widehat z_u^{[v]}&=z_{u-1}^{[v]}+\mu_u^{[v]}.
 \label{eq:report-screen-prediction}
\end{align}
Flag unusual movement when
\begin{equation}
 |z_u^{[v]}-\widehat z_u^{[v]}|>\max\{\log2,4s_u^{[v]}\}
 \quad\text{and}\quad
 |y_u^{[v]}-(e^{\widehat z_u^{[v]}}-1)|\ge10.
 \label{eq:report-movement-flag}
\end{equation}
With insufficient history, the report is unassessed. A separate
revision flag requires successive valid stored values with a $(y+1)$
ratio exceeding two in either direction and a count difference of at
least ten. Flags identify unusual movements or disagreements, not proven
errors; they do not change scoring truth.

\subsection{Feedback}
\label{sec:feedback-designs}
\label{sec:delayed-feedback}
Feedback follows the report and cutoff conventions in
\hyperref[sec:notation]{Notation}, independently of the recalibration method.

\paragraph{Stateful learning.}
Wait $d$ calendar weeks after the first valid report, then consume the
newest report when valid at the next retained forecast cutoff with $t>s$:
\begin{align}
 m_s^{(d)}&=\nu_s+7d\text{ days},
 \label{eq:feedback-maturity}\\
 \tau_d(s)&=\min\{t>s:\iota_t\ge m_s^{(d)},\ y_s^{[\iota_t]}\text{ valid}\},
 &\mathcal R_t^{(d)}&=\{s:\tau_d(s)=t\}.
 \label{eq:feedback-release-set}
\end{align}
An invalid newest report postpones learning; an older version is never
substituted. Use the report at actual consumption time $\iota_t$, even
if a missing forecast delayed it beyond maturity. Each row is learned
at most once; subsequent revisions are ignored.
Plain MultiQT on counts starts at $\tilde\theta_0=0$ and, once active,
updates against previously returned forecasts before computing $q_t$:
\begin{align}
 \tilde\theta_t^\alpha
 &=\tilde\theta_{t-1}^\alpha-\eta_t
   \sum_{s\in\mathcal R_t^{(d)}}
   \bigl(\ind\{y_s^{[\iota_t]}\le q_s^\alpha\}-\alpha\bigr),
 \label{eq:delayed-standard-update}\\
 \theta_t&=\Pi_{\K-b_t}(\tilde\theta_t),\qquad q_t=b_t+\theta_t.
 \label{eq:delayed-standard-forecast}
\end{align}
The hidden state is unprojected; an empty release set gives no update.
The \texttt{log1p} tracker applies the same rule in $\psi$ coordinates.

\paragraph{Weekly re-run.}
At reference week $t$, reconstruct plain MultiQT from zero using reports
current at $\iota_t$, retaining the release positions $j=\tau_d(s)$ and
delay $d$, but omitting invalid current reports:
\begin{align}
 \tilde\theta_j^{\alpha,[t]}
 &=\tilde\theta_{j-1}^{\alpha,[t]}-\eta_j^{[t]}
   \sum_{\substack{s\in\mathcal R_j^{(d)}\\y_s^{[\iota_t]}\text{ valid}}}
   \bigl(\ind\{y_s^{[\iota_t]}\le q_s^{\alpha,[t]}\}-\alpha\bigr),
 \label{eq:rerun-standard-update}\\
 \theta_j^{[t]}&=\Pi_{\K-b_j}(\tilde\theta_j^{[t]}),\qquad
 q_j^{[t]}=b_j+\theta_j^{[t]}.
 \label{eq:rerun-standard-forecast}
\end{align}
Only $q_t^{[t]}$ is returned; earlier returned forecasts are unchanged.
All re-runs keep the stateful $S_0$ and $L_0$ fixed
(Section~\ref{sec:tuning-initialization}). Studentized methods use
\eqref{eq:student-projection}, reconstruct their scale history and return
the base until a required scale exists.
An invalid revised report removes an update, not an expert: the declared
score bank remains fixed (Table~\ref{tab:score-expert-members}).

\paragraph{Experimental withholding.}
The separate screen experiment retains $m_s^{(d)}$ but waits until the
newest eligible report clears the screen. Movement flags are reassessed
at every cutoff; a fresh large-revision flag blocks only the first cutoff
at which that version is current, since unchanged values are not stored
as new versions. Unassessed reports pass. Dated manual reviews could
override flags, but none were supplied. Unresolved flags can prevent a
row from being learned; older reports are never substituted.
Only stateful delays 0 and 1 were tested, excluding re-runs and earlier
A--E mixtures. Tuned parameters stay fixed, but $S_0$ and $L_0$ are
recomputed from the changed release schedule: the comparison combines
postponed updates with changed initialization.

\subsection{Tuning}
\label{sec:tuning-initialization}
Reference scales use the first 20 consumed rows $I_{20}^{(d)}$, which
need not span 20 calendar weeks. Parameter selection uses the tuning
period; these are separate operations.

\paragraph{Reference scales and retrospective initialization.}
The initial error scale pools absolute base errors in tracker units:
\begin{equation}
 S_0=\max\!\left\{\varepsilon_S,
 Q_{0.9}\!\left(\left\{
 |\psi(b_s^\alpha)-\psi(\bar y_s^{(d)})|:
 s\in I_{20}^{(d)},\ \alpha\in\A
 \right\}\right)\right\},\qquad\varepsilon_S=10^{-6}.
 \label{eq:initial-error-scale}
\end{equation}
It sets rates $\eta=\kappa S_0$ and studentized scale bounds
$[.05S_0,S_0]$; studentized step sizes are selected in standardized units.
The reference loss $L_0$ is defined in
Section~\ref{sec:expert-loss-normalization}.

With the default \texttt{replay\_warmup=TRUE}, both constants use the
eventual first 20 releases and apply from the start. Forecasts before
the 20th release therefore use retrospective scale information.
In the active data these rows have tuning target dates, but some reports
arrive after early evaluation cutoffs: tuning-period membership does
not imply availability before evaluation. With replay disabled, learning
waits for warmup; fewer than 20 releases leaves learning inactive.
The constants stay fixed across weekly re-runs; states, weights and
time-varying scales are reconstructed. Evaluation never resets the
tracker, so its initial evaluated correction is generally nonzero.

\paragraph{Cases and selection.}
Tune COVID before flu, and \texttt{log1p} before \texttt{identity}
within each disease. For each pair, run stateful delay 0, delay 1 and
weekly re-run at delay 0. These 12 configurations have two
stages each (24 total). Each starts from defaults; configurations without
tuning rows stop. The hub design tunes horizons $-1,\ldots,3$; the
seasons design tunes $0,\ldots,3$, with $-1$ reusing their choices.

Candidates replay the full history with common $S_0,L_0$ within each
case, tracker and feedback design. Only tuning target weeks are scored, including initialization
rows when valid. Scores use latest archived truth, which need not have
been available before evaluation. Pool case-level means equally over
non-US locations and horizons. The criterion is log quantile loss for
\texttt{log1p} and quantile loss per $100{,}000$ for \texttt{identity},
without $L_0$ normalization.

Stage 1 selects tracker rates separately for \texttt{fixed},
\texttt{leaky}, \texttt{absolute\_slope}, \texttt{studentized},
\texttt{studentized\_base}, \texttt{median\_constant},
\texttt{studentized\_full} and \texttt{studentized\_level\_rms}.
Retain candidates satisfying
\[
 \overline{\mathcal L}_j(\vartheta)
 \le\min_{\vartheta'}\overline{\mathcal L}_j(\vartheta')
       +.01\,\overline{\mathcal L}_0,
\]
then select the smallest mean calibration error. This elbow rule allows
one percentage point above the minimum on the loss-change plot; it does
not estimate curvature.
Stage 2 holds those rates fixed and minimizes loss over $H$ separately
for \texttt{weighted}, \texttt{weighted\_common}, \texttt{shared},
\texttt{shared\_common}, \texttt{restart\_weighted},
\texttt{restart\_grid} and \texttt{score\_learned}.
Both stages break ties by the smaller parameter value. This single pass
avoids a Cartesian search; earlier rate choices are not revisited and
need not give the best joint combination.

\begin{center}\small
\begin{tabularx}{\textwidth}{@{}rlX@{}}
\toprule
Stage & Setting & Candidates\\
\midrule
1 & \texttt{fixed\_multiplier} & .001, .0025, .005, .01, .02, .05, .1, .2, .4\\
1 & \texttt{slope\_min\_multiplier} & .002, .004, .01, .02, .05, .1\\
1 & \texttt{median\_rate} & .003, .01, .03, .1, .3, 1\\
1 & \texttt{leaky\_multiplier} & .001, .002, .004, .008, .02, .05, .1, .2\\
2 & \texttt{H} & 2, 4, 8, 16, 32\\
\bottomrule
\end{tabularx}
\end{center}
Multipliers scale $S_0$; the slope multiplier controls its rate floor.
\texttt{median\_rate} is standardized. $H$ counts processed weight
updates and also sets sharing $\gamma=1/(2H)$ where enabled.

\paragraph{Fixed and inherited settings.}
Learned mixtures use adaptive $\lambda$ from $\lambda_{\mathrm{start}}=2.72$;
$\lambda=0$ gives equal weights in \texttt{exp\_weights()}.
Recent-start count is 8; outer and Decline rates are $.001S_0$ and
$.05S_0$; target-signal threshold/width are 0/.05. Other constants
remain as specified with each method.
\texttt{studentized\_zero} and \texttt{median\_constant\_zero}
inherit their corresponding normal-start rates; \texttt{score\_gated}
inherits $H$ from \texttt{score\_learned}, and outer variants retain
the selected settings of their inner mixtures.

\paragraph{Tuning plots.}
The current \path{multiqt_tuning.pdf} pairs calibration error with
percentage loss change $100(\overline{\mathcal L}_j/\overline{\mathcal L}_0-1)$.
Solid/dashed lines use the different log/identity criteria. Filled rate
points meet the tolerance; open points exceed it; large circles mark
selected values. A horizontal CE line marks the original.
Pooled pages by disease and feedback design precede pages for California,
Texas, New York and Pennsylvania, when available. State pages average
over horizons but circle the pooled selection. $H$ curves condition on
selected rates. The \texttt{on\_grid\_edge} flag marks an endpoint
selection, not evidence of an interior optimum. Embedded plots at
page~\pageref{plot:tuning} retain their recorded earlier procedure.

\subsection{Evaluation}
\label{sec:evaluation-design}
\paragraph{Peaks and evaluation windows.}
\texttt{find\_peaks()} removes nonfinite counts and the incomplete-reporting
interval, then smooths latest counts:
\[
 z_i=(x_{i-1}+x_i+x_{i+1})/3,\qquad
 P_i=z_i-\max\!\left\{\min_{l_i\le j\le i}z_j,
                         \min_{i\le j\le r_i}z_j\right\}.
\]
Endpoints retain the unsmoothed value. Candidates are maxima within
eight retained observations on each side, earliest at ties, satisfying
\[
 P_i>0,\qquad P_i\ge.30z_i,\qquad P_i\ge.10X.
\]
Place each peak at the largest raw count within one retained observation
of its candidate, earliest at ties; keep dates within the forecast target
range. These windows count observations, not calendar weeks across gaps.
Peaks are hindsight labels, used for prediction only by the target oracle.

Assign each target to its nearest peak (earlier at ties), then score
\begin{align}
 W_{\mathrm{overall}}&=E,\notag\\
 W_{\mathrm{before}}&=\{t\in E:-7\le u_t^{\mathrm{peak}}\le0\},\notag\\
 W_{\mathrm{after}}&=\{t\in E:1\le u_t^{\mathrm{peak}}\le8\}.
 \label{eq:evaluation-windows}
\end{align}
Rows count once; the peak week belongs to ``up to peak''. Overall also
includes weeks outside these windows. Missing rows and period boundaries
shorten windows without discarding peaks. Period membership follows the
target date; its nearest peak may lie in another period. Seasonal scores
restrict these sets; the main plots combine both evaluation seasons.

\paragraph{Loss and calibration error.}
For each case, average over selected rows and all 23 quantile levels:
\begin{align}
 \mathcal L^{\mathrm{nat}}(W)
 &=\frac1{mN}\sum_{t\in W}\sum_{\alpha\in\A}
 \rho_\alpha(q_t^\alpha/c_{\mathrm{scale}},y_t^\star/c_{\mathrm{scale}}),
 \label{eq:evaluation-natural}\\
 \mathcal L^{\log}(W)
 &=\frac1{mN}\sum_{t\in W}\sum_{\alpha\in\A}
 \rho_\alpha\!\left(\log(1+\max\{q_t^\alpha,0\}),\log(1+y_t^\star)\right),
 \label{eq:evaluation-log}\\
 \mathrm{CE}(W)
 &=\frac1m\sum_{\alpha\in\A}
 \left|\frac1N\sum_{t\in W}\ind\{y_t^\star\le q_t^\alpha\}-\alpha\right|.
 \label{eq:evaluation-ce}
\end{align}
Columns are \texttt{quantile\_loss}, \texttt{log\_quantile\_loss} and
\texttt{coverage\_error}; tuning uses \texttt{loss}, \texttt{log\_loss}
and \texttt{coverage\_error}. CE averages weeks before taking the absolute
residual. All feedback designs use the same latest truth; empty windows
give missing scores. Evaluation losses are never divided by $L_0$.
Scripts 04/06 display the criterion of the tracker; script 02 displays
loss per $100{,}000$.

\paragraph{Quantile-level diagnostics.}
\path{multiqt_by_level.pdf} repeats the fixed-rate trade-off at each level:
\begin{align}
 \mathcal L_\alpha^{\mathrm{nat}}(W)
 &=\frac1N\sum_{t\in W}
 \rho_\alpha(q_t^\alpha/c_{\mathrm{scale}},y_t^\star/c_{\mathrm{scale}}),\notag\\
 \mathrm{CE}_\alpha(W)
 &=\left|\frac1N\sum_{t\in W}\ind\{y_t^\star\le q_t^\alpha\}-\alpha\right|.
 \label{eq:evaluation-by-level}
\end{align}
Pages fix disease and level, with location rows, horizon columns and
phase symbols. There are 23 COVID pages then 23 flu pages, in increasing
$\alpha$ order; $.1,.5,.9$ are loaded at page~\pageref{plot:by-level}.

\paragraph{All-location trade-offs.}
In \path{multiqt_peak_reasoning_tradeoff.pdf}, columns are horizons
$-1,\ldots,3$; blocks are delay 0, delay 1 and re-run, each showing
overall, up to peak and after peak. Lower left is better. The first
page per disease shows individual non-US cases and large method means
\[
 \left(J^{-1}\sum_c\mathcal L_{j,c},\ J^{-1}\sum_c\mathrm{CE}_{j,c}\right).
\]
These average location scores, not forecasts or pooled coverage
indicators. The first row of \path{multiqt_error_tradeoff.pdf} uses
the same aggregation for fixed rates. Feedback designs have separately
selected parameters, so their differences include retuning.
Axes vary by panel; all-location limits use the 1st--99th percentiles
and mean points, with loss capped near four times the base. Extremes
can be hidden or marked \texttt{>}. Original/common-weight/outer methods
use open circles/triangles/squares or diamonds; exact forecast duplicates
are hidden.

\paragraph{Evaluation tree.}
Script 06 splits pooled disease cases by feedback design, then horizon.
For each window it reports
\begin{align}
 R_j(\mathcal C)
 &=100\left[1-\frac{\sum_{c\in\mathcal C}N_c\mathcal L_{j,c}}
                        {\sum_{c\in\mathcal C}N_c\mathcal L_{0,c}}\right],
 \label{eq:tree-reduction}\\
 P_j(\mathcal C)
 &=J^{-1}\sum_{c\in\mathcal C}\ind\{\mathcal L_{j,c}<\mathcal L_{0,c}\}.
 \label{eq:tree-improvement-share}
\end{align}
\texttt{reduction} is a ratio of week-weighted losses (positive is better);
\texttt{share\_improved} weights cases equally, counts ties as not improved, and is
printed for at least five cases. CE panels also average case scores
with week weights. Each page displays the ten methods with lowest
pooled after-peak loss: a descriptive evaluation ranking, not tuning.
Overlapping horizons/designs are not independent replicates. Tree panels
share axes; clipping is marked \texttt{<} or \texttt{>}.
By contrast, script 02 phase-rate plots show loss change
$100(\mathcal L_j/\mathcal L_0-1)$, where negative is better.

\paragraph{Forecast curves.}
\path{multiqt_method_forecasts.pdf} shows levels $.1,\ldots,.9$ and
weekly loss per $100{,}000$ at a fixed horizon; gray marks tuning,
dotted lines peaks. \path{multiqt_season_forecasts.pdf} adds coverage
and loss trajectories for selected fixed rates.
\path{multiqt_base_forecast_paths.pdf} instead shows each submission
across horizons $-1$ to 3, with later reference dates darker. Its flu
winter 2024/25 peak is within tuning, not held-out evidence.
Weight timelines and archived A--E soft-switch timelines explain
mixtures, rather than providing additional performance scores.
Gap dates are listed in Section~\ref{sec:data-issues}.

\clearpage

\subsection{Methods}
\label{sec:methods}
The methods below use the preceding data, feedback, tuning and evaluation
conventions unless stated otherwise.

\subsubsection{Single trackers}
\label{sec:single-trackers}
These methods maintain one MultiQT tracker for each case.

\paragraph{Fixed rate: \texttt{fixed}.}
Use \eqref{eq:delayed-standard-update}--\eqref{eq:delayed-standard-forecast}
with $\eta=\kappa S_0$ in tracker units. Tune $\kappa$ separately
(default $.05$). Corrections persist across rounds.

\paragraph{Leaky tracker: \texttt{leaky}.}
Add a common decaying shift to persistent MultiQT offsets, as in
\eqref{eq:leaky-shift}--\eqref{eq:identity-leak}. The separately tuned
method also belongs to the score bank.

\paragraph{Slope-dependent rate: \texttt{absolute\_slope}.}
Use the delayed MultiQT update with rate
\begin{equation}
 \eta_t=\min\!\left\{
 \kappa_{\max}S_0,\ \kappa_{\min}S_0+
 \left|\frac{\beta_t^{.5}-\beta_{t-1}^{.5}}{\Delta_t}\right|
 \right\}.
 \label{eq:slope-rate}
\end{equation}
The first slope is zero. Tune the floor multiplier (default $.004$)
and keep the cap at $.08$. The slope compares base medians at the same
horizon across reference dates, per elapsed calendar week; it uses no
future outcomes. Unlike studentization, it changes only new increments.

\paragraph{Persistent restart control: \texttt{restart\_persistent}.}
Keep the first eligible weekly-start expert with $\kappa=.05$
(Section~\ref{sec:expert-combinations}). It learns only from forecasts
made since its start. Its rate is not separately tuned, so this control
for restart mixtures need not reproduce \texttt{fixed}.

\subsubsection{Studentized scores}
\label{sec:studentized-code-updates}
\paragraph{Studentized variants and median initialization.}
In tracker units, use the full base or its median as the centre. After
processing available feedback, return
\begin{equation}
 q_t=\psi^{-1}\!\left[
   \Pi_\K^{A_t^{-2}}\bigl(a_t+A_t\tilde u_t\bigr)
 \right].
 \label{eq:studentized-implementation}
\end{equation}
For a common scale, the projection is the ordinary isotonic projection.
After warmup, the delayed update is applied to the hidden state,
\begin{equation}
 \tilde u_t^\alpha
 =\tilde u_{t-1}^\alpha+
 \eta\sum_{s\in\mathcal R_t^{(d)}}
 \left(\alpha-\ind\{\bar y_s^{(d)}\le q_s^\alpha\}\right).
 \label{eq:studentized-delayed-code}
\end{equation}
Compare each report with its historical forecast $q_s$, reconstructed
inside a re-run when applicable. Project only the forecast; the hidden
state persists across peaks and seasons. The standardized step has
default $\eta=.1$ and is tuned for each normal-start or full-base method.
Zero-start variants inherit the corresponding normal-start rate.

\paragraph{Scales from recent errors or the current forecast.}
The residual window contains five consumed reports, not necessarily
five consecutive calendar weeks. The default scales are
\begin{align}
 \widehat\sigma_{t,\mathrm{med}}
 &=\left[\frac{1}{|\mathcal H_t^{(d)}|}
          \sum_{s\in\mathcal H_t^{(d)}}(e_s^{.5})^2\right]^{1/2},
 \label{eq:studentized-median-rms}\\
 \widehat\sigma_{t,\mathrm{res}}^\alpha
 &=\left[\frac{1}{|\mathcal H_t^{(d)}|}
          \sum_{s\in\mathcal H_t^{(d)}}(e_s^\alpha)^2\right]^{1/2}.
 \label{eq:studentized-level-rms}
\end{align}
Both require five consumed rows. RMS includes bias because the errors
are not centred. The optional robust alternative is
$\operatorname{median}|e_s^\alpha|/\Phi^{-1}(.75)$.

The common scale inferred from the current base is
\begin{equation}
 \widehat\sigma_{t,\mathrm{base}}
 =\frac{\beta_t^{.9}-\beta_t^{.1}}
        {\Phi^{-1}(.9)-\Phi^{-1}(.1)}.
 \label{eq:studentized-base-width}
\end{equation}
The denominator is approximately $2.563$. The earlier
\texttt{studentized\_levels} method used the level-specific base scales
\begin{equation}
 \widehat\sigma_{t,\mathrm{pair}}^\alpha
 =\frac{\beta_t^{1-a_\alpha}-\beta_t^{a_\alpha}}
        {\Phi^{-1}(1-a_\alpha)-\Phi^{-1}(a_\alpha)}.
 \label{eq:studentized-paired-width}
\end{equation}
This construction requires a quantile grid symmetric about $.5$.
Base-width scales respond immediately to changes in the current forecast;
residual scales respond only after the corresponding reports arrive.

Floor raw scales at $10^{-6}$, then clip to $[.05S_0,S_0]$ once
$S_0$ exists; denote these bounded scales by $\sigma$. The bounds limit
how much the scale suppresses or amplifies accumulated corrections.
They are implementation defaults, not requirements of studentization.

\begingroup
\small
\setlength{\tabcolsep}{4pt}
\renewcommand{\arraystretch}{1.18}
\begin{longtable}{@{}>{\raggedright\arraybackslash}p{.30\textwidth}
 >{\raggedright\arraybackslash}p{.23\textwidth}
 >{\raggedright\arraybackslash}p{.38\textwidth}@{}}
\caption{Studentized methods. Scales refer to the bounded versions above.
Initial states are specified before the first update.}
\label{tab:studentized-implementations}\\
\toprule
Method ID & Centre and scale & Initialization and purpose\\
\midrule
\endfirsthead
\multicolumn{3}{l}{\tablename\ \thetable\ (continued)}\\
\toprule
Method ID & Centre and scale & Initialization and purpose\\
\midrule
\endhead
\midrule
\multicolumn{3}{r}{Continued on next page}\\
\endfoot
\bottomrule
\endlastfoot
\texttt{studentized}
 & $a_t^\alpha=\beta_t^{.5}$; common $\sigma_{t,\mathrm{med}}$
 & $\tilde u_{\mathrm{init}}^\alpha=\Phi^{-1}(\alpha)$.
   Reconstruct the distribution around the base median using recent errors.\\
\texttt{studentized\_base}
 & $a_t^\alpha=\beta_t^{.5}$; common $\sigma_{t,\mathrm{base}}$
 & Normal initial state. Reconstruct the distribution from the median
   and the central width of the current base forecast.\\
\texttt{median\_constant}
 & $a_t^\alpha=\beta_t^{.5}$; common constant $\sigma$
 & Normal initial state; freeze $\sigma_{t,\mathrm{med}}$ at the first
   reference week after both required warmups. Separates replacing the base by its median from
   subsequently changing the scale.\\
\texttt{studentized\_zero}
 & $a_t^\alpha=\beta_t^{.5}$; common $\sigma_{t,\mathrm{med}}$
 & $\tilde u_{\mathrm{init}}=0$. Start all levels at the median;
   use the rate selected for \texttt{studentized}.\\
\texttt{median\_constant\_zero}
 & $a_t^\alpha=\beta_t^{.5}$; common constant $\sigma$
 & Zero initial state; use the scale and rate of
   \texttt{median\_constant}.\\
\texttt{studentized\_full}
 & $a_t^\alpha=\beta_t^\alpha$; common $\sigma_{t,\mathrm{base}}$
 & $\tilde u_{\mathrm{init}}=0$. Keep the full base and rescale its
   accumulated corrections together.\\
\texttt{studentized\_level\_rms}
 & $a_t^\alpha=\beta_t^\alpha$; $\sigma_{t,\mathrm{res}}^\alpha$
 & Zero initial state. Set the scale at each level from recent errors at that level;
   weighted isotonic projection.\\
\end{longtable}
\endgroup

The paired-width method \texttt{studentized\_levels}, removed on
11 November 2026, is retained here only to interpret earlier plots.
The scale of \texttt{median\_constant} is fixed within each reconstructed
history but can change between re-runs; $S_0$ and $L_0$ cannot.
Zero-start variants compare initialization at the same selected rate.
Their first active forecast may already include a feedback update.

The residual-median methods also allow $t_5^{-1}(\alpha)$ initialization;
the base-width median method always uses normal quantiles. It retains
the median and bounded central width but imposes a symmetric normal
shape in tracker units, so it need not reproduce the original forecast.
Until a required residual scale exists, return the original full base.
Feedback consumed while waiting is not replayed; later feedback on
those rounds is compared with the original forecasts.

Studentization rescales the accumulated correction, rather than only
the newest increment. For full-base \texttt{log1p} methods without crossings,
\begin{equation}
 q_t^\alpha=(1+b_t^\alpha)\exp(\sigma_t^\alpha\tilde u_t^\alpha)-1,
 \label{eq:log-studentized-nonadditive}
\end{equation}
so the count correction is non-additive. With \texttt{identity}, it is
additive with a changing multiplier. Quantile-specific scales retain the
coverage target, but no calibration guarantee is established here for
those variants.

\subsubsection{Weighted experts}
\setcounter{secnumdepth}{4}
\label{sec:expert-weighting}
An expert is a complete forecasting procedure; a bank is the set being
combined. Weights learn from quantile losses separately for each
location--horizon case. A fixed learning rate still changes the weights;
an adaptive rate changes their responsiveness. The formulas first
assume a fixed bank with all experts eligible.

\paragraph{Mixtures and feedback losses}
With weights specific to each quantile level, compute
\begin{equation}
 \hat q^\alpha=\sum_{k=1}^K w_k^\alpha f_k^\alpha,
 \qquad w_k^\alpha\ge0,\quad \sum_{k=1}^K w_k^\alpha=1,
 \qquad q=\Pi_\K(\hat q).
 \label{eq:expert-mixture}
\end{equation}
Per-quantile weights can introduce crossings, repaired by isotonic
projection. Common weights replace $w_k^\alpha$ by $w_k$ and preserve
order when every expert is ordered. All mixtures average counts,
including for \texttt{log1p} trackers.

On receiving feedback, score the expert forecasts for the released row:
\begin{align}
 L_{k,n}^\alpha
 &=\rho_\alpha\!\left(f_{k,s(n)}^\alpha/c_{\mathrm{scale}},
                      Y_n/c_{\mathrm{scale}}\right),
 &L_{k,n}&=\frac{1}{m}\sum_{\alpha\in\A}L_{k,n}^\alpha,
 \label{eq:expert-feedback-loss}\\
 \ell_{k,n}^\alpha&=L_{k,n}^\alpha/L_0,
 &\ell_{k,n}&=\frac{1}{m}\sum_{\alpha\in\A}\ell_{k,n}^\alpha.
 \label{eq:normalized-expert-loss}
\end{align}
The same $L_0$ applies to all experts and levels within a case and
feedback design (Section~\ref{sec:expert-loss-normalization}). Score the
originally computed forecasts in stateful learning and the reconstructed
forecasts in a re-run.

\paragraph{Exponential updates and optional sharing}
Before optional sharing, update separately at each level:
\begin{equation}
 \bar w_{k,n+1}^\alpha
 =\frac{w_{k,n}^\alpha\exp(-\lambda_n^\alpha\ell_{k,n}^\alpha)}
        {\sum_{j=1}^K w_{j,n}^\alpha
         \exp(-\lambda_n^\alpha\ell_{j,n}^\alpha)}.
 \label{eq:level-expert-weight-step}
\end{equation}
With common weights, there is one update using the mean loss:
\begin{equation}
 \bar w_{k,n+1}
 =\frac{w_{k,n}\exp(-\lambda_n\ell_{k,n})}
        {\sum_{j=1}^K w_{j,n}\exp(-\lambda_n\ell_{j,n})}.
 \label{eq:normalized-expert-weight-step}
\end{equation}
Lower loss increases relative weight; larger $\lambda$ strengthens the
response. Weights start uniformly, and $\lambda=0$ keeps equal weights.

When sharing is enabled, follow the exponential step by
\begin{equation}
 w_{k,n+1}^\alpha
 =(1-\gamma)\bar w_{k,n+1}^\alpha+\gamma/K.
 \label{eq:expert-fixed-share}
\end{equation}
Omit $\alpha$ for common weights. Sharing helps downweighted experts
recover: use $\gamma=1/(2H)$ when enabled and zero otherwise.
Table~\ref{tab:expert-combinations} distinguishes sharing from common
weights across levels.

\paragraph{Adaptive learning rates}
\label{sec:adaptive-expert-weighting}
The first $H$ processed updates use $\lambda_{\mathrm{start}}$.
Record the weighted variance of each update using its weights before the
exponential step and sharing:
\begin{align}
 M_n^\alpha
 &=\sum_{k=1}^K w_{k,n}^\alpha
   \left(\ell_{k,n}^\alpha-
         \sum_{j=1}^K w_{j,n}^\alpha\ell_{j,n}^\alpha\right)^2,\notag\\
 M_n
 &=\sum_{k=1}^K w_{k,n}
   \left(\ell_{k,n}-\sum_{j=1}^K w_{j,n}\ell_{j,n}\right)^2.
 \label{eq:expert-loss-moments}
\end{align}
Treat variances at most $10^{-12}\sum_k w_{k,n}^\alpha
(\ell_{k,n}^\alpha)^2$ as zero, and analogously for common weights,
to avoid large rates from numerical rounding.
For per-quantile weighting, the rate is
\begin{equation}
 \lambda_n^\alpha=
 \begin{cases}
  \lambda_{\mathrm{start}},&1\le n\le H,\\[3pt]
  \displaystyle\sqrt{\frac{\log(2KH)+1}
                    {\sum_{r=n-H}^{n-1}M_r^\alpha}},&n>H.
 \end{cases}
 \label{eq:adaptive-expert-rate}
\end{equation}
Only the previous $H$ variances enter; a zero sum retains the previous
finite rate. Greater loss differences reduce the rate, whereas a large
error shared by all experts does not increase the variance.

Per-quantile weighting maintains a vector, variance history and rate
for each level. Common weighting uses one of each: average losses
across levels \emph{before} computing $M_r$, then omit $\alpha$ in
\eqref{eq:adaptive-expert-rate}. This is not an average of separately
learned weights or per-quantile variances.

\subparagraph{Later starts and implementation.}
Only restart experts that forecast the released row are eligible.
Normalize their weights for the update, then restore their total mass;
other weights stay unchanged. Use the eligible count in the rate formula.
A new expert receives $1/K$ of the enlarged bank, preserving existing
relative weights. Updates with fewer than two eligible experts do not
count toward $H$.

\paragraph{Expert loss normalization}
\label{sec:expert-loss-normalization}
Using the quantile loss in~\eqref{eq:quantile-loss}, define one reference loss
per case and delay:
\begin{equation}
 L_0=\max\!\left\{\varepsilon_L,
 \frac{1}{20|\A|}\sum_{s\in I_{20}^{(d)}}\sum_{\alpha\in\A}
 \rho_\alpha\!\left(\frac{b_s^\alpha}{c_{\mathrm{scale}}},
                    \frac{\bar y_s^{(d)}}{c_{\mathrm{scale}}}\right)\right\},
 \qquad\varepsilon_L=10^{-6}.
 \label{eq:expert-reference-loss}
\end{equation}
Losses are relative to the initial mean base loss: $\ell=.8$ means
80\% of that reference, not necessarily a 20\% improvement in the current
week. Values can exceed one. Weighting uses scaled counts even for
\texttt{log1p} trackers.

\subparagraph{Initial and fixed-rate weight updates.}
Dividing by $L_0$ preserves rankings and makes weight changes depend on
relative loss gaps. A change of count units scales both $L$ and $L_0$
equally, leaving weights unchanged when the numerical floor is inactive.
This matters throughout fixed-rate runs and during the first $H$
adaptive-run updates at $\lambda_{\mathrm{start}}$; these are distinct
from the 20 rows used to estimate $L_0$.

For equal initial weights, losses $80$ and $100$, $L_0=100$ and
$\lambda_{\mathrm{start}}=1$ give weights about $.55$ and $.45$ before
sharing. The same numerical rate on raw losses would put almost all
weight on the first expert.

\subparagraph{Effect on adaptive learning rates.}
After $H$ updates, the rule in~\eqref{eq:adaptive-expert-rate} already
adapts to the variation between expert losses. Conditional on the same
pre-update weight vectors at updates $n-H,\ldots,n$ and a positive
variance sum, normalization gives, separately at each level,
\begin{align}
 M_r^{\alpha,\mathrm{norm}}
 &=M_r^{\alpha,\mathrm{raw}}/L_0^2,
 &\lambda_n^{\alpha,\mathrm{norm}}
 &=L_0\lambda_n^{\alpha,\mathrm{raw}},\notag\\
 \lambda_n^{\alpha,\mathrm{norm}}\ell_{k,n}^\alpha
 &=\lambda_n^{\alpha,\mathrm{raw}}L_{k,n}^\alpha.
 \label{eq:expert-normalization-cancellation}
\end{align}
The identity also holds for common weights. The numerical zero-variance
threshold scales by $L_0^{-2}$ too. Thus the loss and rate changes cancel
in the exponential update, conditional on the same weight history.
Runs with the same numerical starting rate can still differ: their
first $H$ updates change later weights and variances.

Equivalent runs require matching settings and
$\lambda_{\mathrm{start}}^{\mathrm{raw}}
=\lambda_{\mathrm{start}}^{\mathrm{norm}}/L_0$, including the same scaling
of retained rates when the variance sum is zero. Keep normalization
throughout a run so incoming losses, stored variances and rates remain
in consistent units. After adaptation, $L_0$ is not an additional
reduction in update strength.

\subparagraph{Effect across quantile levels and re-runs.}
A single $L_0$ does not equalize losses across levels. Per-quantile
weights use the loss and adaptive rate of each level; common weights
use mean losses. Re-runs reconstruct losses but retain the stateful
$L_0$. Evaluation losses are not normalized by $L_0$, and normalization
alone guarantees neither accuracy nor calibration.

\paragraph{Expert combinations implemented in the code}
\label{sec:expert-combinations}
\label{sec:phase-specific-experts}
\label{sec:shared-quantile-weights}
Table~\ref{tab:expert-combinations} groups the current combinations.
Every learned mixture uses the adaptive variance-based rate above.

\begin{center}\small
\renewcommand{\arraystretch}{1.2}
\begin{tabularx}{\textwidth}{@{}>{\raggedright\arraybackslash}p{.32\textwidth}X@{}}
\toprule
Method IDs & Weighting and bank\\
\midrule
\texttt{uniform} & Equal weights on 11 persistent rate experts.\\
\texttt{weighted}, \texttt{shared} & Same rate bank; per-quantile weights, without/with sharing.\\
\texttt{weighted\_common}, \texttt{shared\_common} & Same rate bank; common weights, without/with sharing.\\
\texttt{restart\_uniform} & Equal weights on all weekly starts, rate $.05S_0$.\\
\texttt{restart\_weighted} & Same starts; adaptive $\lambda$, sharing, per-quantile weights.\\
\texttt{restart\_grid} & Weekly starts at $.01S_0,.05S_0,.2S_0$; adaptive $\lambda$, sharing, per-quantile weights.\\
\texttt{restart\_recent} & Equal weights on the eight most recent starts, rate $.05S_0$.\\
\texttt{score\_equal} & Equal weights on the declared score bank.\\
\texttt{score\_learned} & Same score bank; adaptive $\lambda$, sharing, per-quantile weights.\\
\texttt{score\_gated} & Mix \texttt{score\_learned} with the original using the target-specific gate $G_{t,h}^{T}$.\\
\bottomrule
\end{tabularx}
\end{center}
\refstepcounter{table}\label{tab:expert-combinations}
\begin{center}\small Table \thetable: Current expert combinations.
\end{center}

The five \texttt{\_outer} variants of \texttt{shared},
\texttt{shared\_common}, \texttt{restart\_uniform},
\texttt{restart\_weighted} and \texttt{restart\_grid} add persistent
MultiQT at $.001S_0$ to their inner mixture, retaining its settings
and learning from their own final forecasts. Rate and $H$ selection
follow Section~\ref{sec:tuning-initialization}.

\subparagraph{Rate and weekly-start banks.}
The rate bank uses
\[\kappa\in\{0,.01,.015,.02,.03,.05,.075,.1,.15,.2,.3\};\]
all experts start together and the zero-rate member is the original.
Weekly-start banks add a fresh copy at every retained reference week
when initialization permits; with warmup replay, starts begin in the
first round. Each copy starts with zero correction and learns only
from feedback on forecasts made since that start. Older copies keep
their accumulated corrections. There is no permanent zero-rate expert.

For \texttt{restart\_recent}, $\mathcal B_t$ contains the most recent
$K_{\mathrm{recent}}=8$ copies, or all copies if fewer have started. Then
\begin{equation}
 q_t^\alpha=|\mathcal B_t|^{-1}\sum_{k\in\mathcal B_t}f_{k,t}^\alpha.
 \label{eq:recent-start-mixture}
\end{equation}
This controls the age of corrections without detecting peaks. Count
starts, not elapsed weeks across gaps. A new expert receives weight
before feedback, so a new color band alone does not show lower loss.

\subparagraph{The six score-based experts and fixed membership.}
\label{sec:score-expert-bank}
The bank compares ways to retain or rebuild the base and correct its
errors. For the scores below, $r=\psi(y)$; project and invert as in
\eqref{eq:studentized-implementation}. Fixed, leaky and studentized
experts use separately selected rates; Decline uses $.05S_0$ and the
original does not update. This compares complete procedures, not
scores at a common rate.

Membership is fixed across reference weeks and re-runs, even when
forecasts coincide. Exclude Decline at $h=-1$, where it equals the
original by construction. Do not deduplicate complete forecast histories:
later forecasts or revisions must not change earlier membership.

\begingroup\small
\renewcommand{\arraystretch}{1.18}
\begin{longtable}{@{}>{\raggedright\arraybackslash}p{.27\textwidth}
 >{\raggedright\arraybackslash}p{.32\textwidth}
 >{\raggedright\arraybackslash}p{.33\textwidth}@{}}
\caption{The declared score bank. Scales are the bounded scales defined earlier;
Decline is omitted at $h=-1$ under the fixed-membership correction.}
\label{tab:score-expert-members}\\
\toprule
Member & Score or forecast & Purpose\\
\midrule\endfirsthead
\toprule Member & Score or forecast & Purpose\\\midrule\endhead
\bottomrule\endfoot
\texttt{original} & $f_{k,t}^\alpha=b_t^\alpha$ & Preserve the base as an option.\\
\texttt{fixed} & $S_{k,t}^\alpha(y)=r-\beta_t^\alpha$ & Persistent additive correction; $\eta=\kappa S_0$.\\
\texttt{leaky} & Fixed MultiQT plus $v_t$ in \eqref{eq:leaky-shift} & Short-lived common directional correction alongside persistent offsets.\\
\texttt{studentized\_base} & $S_{k,t}^\alpha(y)=(r-\beta_t^{.5})/\sigma_{t,\mathrm{base}}$ & Rebuild around the median with current central width and normal initial thresholds.\\
\texttt{studentized\_full} & $S_{k,t}^\alpha(y)=(r-\beta_t^\alpha)/\sigma_{t,\mathrm{base}}$ & Keep the full base; rescale accumulated corrections with its central width.\\
\texttt{decline} & $S_{k,t}^\alpha(y)=r-\beta_t^\alpha$, updated only for $c_{s,h}^{\mathrm{ctx}}=1$ & Learn and apply a persistent correction in declining forecast context; return the original otherwise.\\
\end{longtable}\endgroup
Initialization and scales follow Table~\ref{tab:studentized-implementations}.
Leaky and Decline change state updates, rather than defining new static
residual scores. \path{studentized_level_rms} remains a standalone
method; \path{studentized_levels} appears only in earlier results.

\subparagraph{The leaky expert.}
Besides its persistent MultiQT offsets, this expert has one common
shift on the log scale, initially zero. After warmup, its update is
\begin{equation}
 v_t=\operatorname{clip}_{[-.2,.2]}\!\left[
  2^{-\Delta_t}v_{t-1}+
  .2\sum_{s\in\mathcal R_t^{(d)}}\frac1m
       \sum_{\alpha\in\A}
       \left(\alpha-\ind\{\bar y_s^{(d)}\le f_{\mathrm{leaky},s}^\alpha\}\right)
 \right],
 \label{eq:leaky-shift}
\end{equation}
The shift halves per calendar week, including gaps. Sum reports within
a batch, then cap the shift. Errors in the same direction reinforce it;
opposite errors across levels can cancel. It therefore targets a common
directional error. Ordinary offsets persist and update at each level.

For \texttt{log1p}, add $v_t$ after projection and before inversion.
For \texttt{identity}, apply it to the ordinary projected count forecast
$q_{\mathrm{plain},t}^\alpha$ as
\begin{equation}
 f_{\mathrm{leaky},t}^\alpha
 =\bigl(1+(q_{\mathrm{plain},t}^\alpha)_+\bigr)e^{v_t}-1+\min(q_{\mathrm{plain},t}^\alpha,0).
 \label{eq:identity-leak}
\end{equation}
This gives the same log-shift for nonnegative counts and extends it to
negative forecasts. These trackers do not clip final forecasts at zero.


\subparagraph{Target-specific signal and Decline.}
The signal compares the target horizon with the largest available
earlier median in the same forecast path:
\begin{align}
 D_{t,h}&=\max_{\substack{r=-1,\ldots,h-1\\Z_{t,r}\ \mathrm{available}}}
               Z_{t,r}-Z_{t,h},\notag\\
 c_{t,h}^{\mathrm{ctx}}&=\ind\{D_{t,h}>\log(1+\delta_G)\},&
 G_{t,h}^{T}&=\left[1+\exp\left\{-\frac{D_{t,h}-\log(1+\delta_G)}
                                      {\log(1+w_T)}\right\}\right]^{-1}.
 \label{eq:target-context-signal}
\end{align}
The target median and at least one earlier median must be available
by the cutoff at the exact reference date, excluding filled forecasts.
Without this information, including at $h=-1$, set the context and
gate to zero. The fixed settings are $\delta_G=0$ and $w_T=.05$.
When $D_{t,h}=0$, the hard context is zero but the gate is $.5$.
The signal describes a decline in the base forecast path; it is not
a calibrated probability of being after the true peak.

Decline starts with zero hidden offsets and uses
\begin{align}
 \tilde\theta_t^\alpha
 &=\tilde\theta_{t-1}^\alpha+.05S_0
   \sum_{s\in\mathcal R_t^{(d)}}c_{s,h}^{\mathrm{ctx}}
   \left(\alpha-\ind\{\bar y_s^{(d)}\le f_{\mathrm{decline},s}^\alpha\}\right),
 \label{eq:decline-tracker-update}\\
 f_{\mathrm{decline},t}
 &=\begin{cases}
 \psi^{-1}\!\left[\Pi_\K(\beta_t+\tilde\theta_t)\right],
       &c_{t,h}^{\mathrm{ctx}}=1,\\
 b_t,&c_{t,h}^{\mathrm{ctx}}=0.
 \end{cases}
 \label{eq:decline-tracker-forecast}
\end{align}
The historical context of row $s$ determines whether its feedback
updates this state; the current context determines whether to apply
the correction. Offsets persist between declining periods. The same
warmup and re-run conventions apply as for Fixed MultiQT.

\subparagraph{Equal, learned and gated score mixtures.}
\texttt{score\_equal} averages the declared experts. For
\texttt{score\_learned}, learn one vector per quantile level using
the normalized losses, adaptive rate and sharing above, then repair
crossings. The gated mixture is
\begin{equation}
 q_{\mathrm{score\_gated},t}^\alpha
 =(1-G_{t,h}^{T})q_{\mathrm{score\_learned},t}^\alpha
   +G_{t,h}^{T}b_t^\alpha.
 \label{eq:score-gated-forecast}
\end{equation}
The gate changes the returned forecast, retaining the weights, $H$
and training losses of \texttt{score\_learned}. It introduces no second
weight vector. A missing signal returns the learned forecast; using
the same gate across levels preserves order.

\subparagraph{Earlier A--E experiments.}
The results in Section~\ref{sec:results} use the earlier seven-member
bank: original, fixed, leaky, studentized base width, studentized full
base width, studentized level RMS and studentized paired level scales.
They predate the six-member bank and the revised tuning procedure.
A (\texttt{score\_single}) learned one common vector. B
(\texttt{score\_switch}) used two softly updated vectors with a gate
from forecast-path, nowcast and forecast-revision signals. C
(\texttt{score\_target\_switch}) used the target-specific gate instead.
D (\texttt{score\_context}) trained each of two vectors only on its
historical context; E (\texttt{score\_fast\_decline}) increased the
learning rate during current declining context. Their hard-selection
and no-sharing variants are also historical comparisons.
These experiments used common weights, $H=4$ and sharing $1/8$
where enabled. Their adaptive rates used uncentered second moments,
and the weekly-start mixtures used a fixed tuned $\lambda$.
Their plots do not establish results for the current bank or for the
new variance-based weighting. The earlier formulas are preserved in
the archived notes accompanying those plots.

\setcounter{secnumdepth}{3}

\subsubsection{Reference benchmarks}
\label{sec:reference-benchmarks}
The references are the original forecast, a hindsight control and two
production configurations of Delphi.

\paragraph{Original forecast: \texttt{original}.}
Return $q_t=b_t$ without recalibration or learning. This is the reference
for the reported loss changes and reductions.

\paragraph{Target oracle: \texttt{oracle\_target}.}
This hindsight control returns
\begin{equation}
 q_t^{\mathrm{oracle}}=
 \begin{cases}
 q_t^{\mathrm{fixed}},&T_t\le p(t),\\
 b_t,&T_t>p(t).
 \end{cases}
 \label{eq:target-oracle}
\end{equation}
This bypasses correction for postpeak targets, including forecasts
made before the peak. Fixed MultiQT still learns from its own forecasts.
The control tests bypassing correction with perfect phase information;
it is not deployable. Setting the rate to zero would instead freeze
the offset, giving a different control.

\paragraph{Delphi benchmarks: separate scales and initialization.}
Both run for both diseases and all feedback designs, independently of
the identity/log tracker setting. States start at zero without $S_0$
or a 20-row warmup; constants are not tuned. In their own units, use
\begin{equation}
 \epsilon_t^\alpha=
 \sum_{s\in\mathcal R_t^{(d)}}
 \left(\alpha-\ind\{\chi(\bar y_s^{(d)})\le\xi_s^\alpha\}\right).
 \label{eq:sqrt-feedback-errors}
\end{equation}
In a weekly re-run, use the reports and reconstructed historical forecasts of that re-run in this expression, as in the feedback section.

\paragraph{Adaptive fast and slow offsets: \texttt{ref\_op}.}
The COVID production configuration uses $\chi(x)=\sqrt{x}$ on raw counts,
with rates
\begin{align}
 \eta_t^F&=\max\!\left\{10^{-3},\ .03Q_{.9}\!\left(
 \left\{|\sqrt{\bar y_s^{(d)}}-\sqrt{b_s^\alpha}|:
 s\in\mathcal F_{t,20}^{(d)},\ \alpha\in\A\right\}\right)\right\},
 \notag\\
 \eta_t^S&=.003Q_{.9}\!\left(
 \left\{|\sqrt{\bar y_s^{(d)}}-\sqrt{b_s^\alpha}|:
 s\in\mathcal F_t^{(d)},\ \alpha\in\A\right\}\right).
 \label{eq:ref-op-rates}
\end{align}
Both are zero before the first release. The fast rate uses up to 20
recent consumed rows; the slow rate uses all consumed rows. Both pool
absolute base errors over levels and are recomputed as feedback arrives.
The hidden states and forecast are
\begin{align}
 \tilde\theta_t^{F,\alpha}
 &=.9\tilde\theta_{t-1}^{F,\alpha}+\eta_t^F\epsilon_t^\alpha,
 &\tilde\theta_t^{S,\alpha}
 &=\tilde\theta_{t-1}^{S,\alpha}+\eta_t^S\epsilon_t^\alpha,
 \notag\\
 \xi_t&=\Pi_\K(\sqrt{b_t}+\tilde\theta_t^F+\tilde\theta_t^S),
 &q_t^\alpha&=(\max\{\xi_t^\alpha,0\})^2.
 \label{eq:ref-op-update}
\end{align}
The fast correction loses 10\% per processed reference week, even without
feedback; gaps do not compound decay. The slow correction persists.
Re-runs reconstruct both rates and states. These rolling scales are
separate from the frozen $S_0$.

\paragraph{Constant square-root rate: \texttt{sqrt\_constant}.}
The flu production benchmark of Delphi uses
$\chi(x)=\sqrt{x/c_{\mathrm{pop}}}$, with
\begin{align}
 \tilde\theta_t^\alpha
 &=\tilde\theta_{t-1}^\alpha+.018\epsilon_t^\alpha,
 \notag\\
 \xi_t&=\Pi_\K\!\left(\sqrt{b_t/c_{\mathrm{pop}}}+\tilde\theta_t\right),
 &q_t^\alpha&=c_{\mathrm{pop}}(\max\{\xi_t^\alpha,0\})^2.
 \label{eq:sqrt-constant-update}
\end{align}
The rate $.018$ is in square-root counts-per-$100{,}000$ units.
There is no adaptive rate, decay or second state. Delphi selected
this constant using flu seasons 2024/25 and 2025/26, which overlap
our evaluation seasons. Although we do not retune it, this is not
an independently selected held-out benchmark for those seasons.

\paragraph{Zero boundary and plots.}
Square-root methods return zero when $\xi_t^\alpha<0$ but learn from
untruncated $\xi_s^\alpha$. At zero outcomes, learning coverage can
differ from coverage of the returned forecast. Both appear under
``Single trackers'' in script 04 and in the evaluation tree of script 06.

\clearpage
\section{Results and where to see them}
\label{sec:results}
\emph{These results describe the archived run of 10--11 November 2026:}
the seven-member score bank and A--E, before the six-member bank,
variance-based weighting, simplified tuning and fixed-membership
correction. The revised methods require new scores and plots.

The run uses \texttt{hub\_first\_months}, protocol
\texttt{latest-at-use-v2}, identity tracking and evaluation target weeks
from June 2025 through September 2026. Unless specified otherwise,
reductions pool 52 non-US locations, horizons $-1,\ldots,3$ and three
feedback designs, with week-weighted pooling as in~\eqref{eq:tree-reduction}.
Positive percentages mean lower quantile loss per $100{,}000$ than the
original. The US is separate; longer-history and \texttt{log1p} evaluation
remain separate work.

\subsection{Do the full quantiles add value beyond the median?}
Retaining the base width matters more than retaining the detailed shape
for COVID; for flu, the full quantile curve performs better. The pooled
overall reductions are:
\begin{center}\small
\begin{tabular}{@{}lrr@{}}
\toprule
Method & COVID & Flu\\
\midrule
Median, constant sigma & $-27.62$ & $-34.31$\\
Studentized median (residual RMS) & $-11.89$ & $-28.13$\\
Studentized median (base width) & $-.61$ & $-9.60$\\
Studentized full base (base width) & $-1.28$ & $-3.68$\\
Studentized full base (level scales) & $-1.10$ & $-4.38$\\
Studentized full base (level RMS) & $.08$ & $-.59$\\
\bottomrule
\end{tabular}
\end{center}
\texttt{studentized\_base} replaces the detailed shape while retaining
the median and CMU central width. The constant-scale and residual-RMS
median methods also replace that width, worsening both diseases.
\texttt{median\_constant} is equivalent to raw median-residual tracking
up to fixed rescaling of its rate and initial offsets. Full-base
studentization retains $\beta_t^\alpha$ and rescales the correction
(Table~\ref{tab:studentized-implementations}). These tuned comparisons
do not support using the median alone, but do not isolate all shape
and width information: initialization and $S_0$ still use the full base.

\emph{Plots:} \path{multiqt_peak_reasoning_tradeoff.pdf}, pp.~1 and~7,
``Overall'' rows, compares loss and CE by horizon and feedback design.
For forecast and weekly-loss curves, use the files in
\path{multiqt_method_forecasts/}: \path{covid_horizon2_delay0.pdf} and
\path{flu_horizon2_delay0.pdf}, pp.~6--7, showing the US median/full-base
groups. These are loaded from notes pages~\pageref{plot:covid-studentized-curves}
and~\pageref{plot:flu-studentized-curves}.

\subsection{What the studentized forecast curves show}
\label{sec:evaluation-forecast-plots}
At $h=2$, delay 0, US COVID full-base curves remain close to the original
and Fixed MultiQT, retaining the lag around the summer 2025 wave.
Constant-scale median reconstruction keeps broad tails during low
incidence, including negative lower and elevated upper quantiles,
with additional loss. \emph{Plots:} COVID forecast PDF, p.~7 (full
base) and p.~6 (median), forecast and weekly-loss rows.

For US flu around January 2026, full-base width and level-scale methods
raise central quantiles and add loss late in the wave; level RMS stays
closer to the original. Median residual RMS produces a pronounced
negative lower tail after the peak. These curves illustrate how a
changing scale can amplify an accumulated correction.
\emph{Plots:} flu forecast PDF, pp.~6--7, January--March 2026 forecast
and weekly-loss rows.

\subsection{COVID, flu, the US and the selected four states}
\label{sec:us-largest-states}
The selected states are California, Texas, New York and Pennsylvania,
not the four largest by population. In the initial rate grid at $h=2$,
delay 0, after peaks, COVID $\kappa=.10$ lowers CE but raises loss in
the US, Texas, New York and Pennsylvania; California worsens on both.
New York improves both at $\kappa=.01$. For flu at $\kappa=.03$, US
and Texas loss is nearly unchanged while CE improves; the other three
states worsen on both. Better coverage therefore need not mean lower
loss. \emph{Plots:} \path{multiqt_error_tradeoff.pdf}, pp.~1 (COVID)
and~4 (flu), $h=2$, US/state rows, after-peak symbols;
\path{multiqt_phase_rate_tradeoffs.pdf}, pp.~2--6, $h=2$, after-peak
curves. Script 02 labels $\kappa$ as $m$.

For tuned methods after peaks, median/base-width reconstruction
improves COVID loss in the US, Texas and New York but worsens California
and Pennsylvania. Full-base level RMS improves flu loss in the US and
Texas but worsens the other three states.
\emph{Plots:} \path{multiqt_peak_reasoning_tradeoff.pdf}, pp.~2--6
(COVID) and~8--12 (flu), after-peak rows across horizons/designs.

Across non-US locations, A--E give only small overall COVID gains and
do not remove the pooled deterioration after peaks. No tested online
method improves pooled flu loss overall. Keeping the eight most recent
starts improves substantially over retaining all starts, with overall
COVID gains in California and Texas but no after-peak gain there.
\emph{Plots:} \path{multiqt_peak_reasoning_tradeoff.pdf}, pp.~1 and~7
(all methods), pp.~3--4 (California/Texas starts);
\path{multiqt_evaluation_tree.pdf}, pp.~1 and~7, root overall/after-peak
panels. At $h=2$, delay 0, COVID forecast pp.~10 and~17 and
\path{multiqt_weight_timelines/covid_horizon2_delay0.pdf}, pp.~6 and~9,
show the corresponding California/Texas curves and weights. The tree
selects ten methods by after-peak loss; use the trade-off plot for the
full comparison.

\subsection{Delayed feedback and weekly re-runs}
\label{sec:feedback-delay-results}
Waiting one additional week can help, but re-running is not consistently
best. For context weighting D, the reductions are:
\begin{center}\small
\begin{tabular}{@{}lccc@{}}
\toprule
 & Delay 0 & Delay 1 & Re-run\\
\midrule
COVID: overall / after peak & $+.71\%/+.65\%$ & $+.13\%/-2.09\%$ & $-.06\%/-2.48\%$\\
Flu: overall / after peak & $-2.02\%/-5.79\%$ & $-1.10\%/-3.94\%$ & $-2.75\%/-7.88\%$\\
\bottomrule
\end{tabular}
\end{center}
Delay 1 trades later learning for a later report, without guaranteeing
better accuracy. Re-runs reconstruct tracker states from the report
versions available then. Reports and tuned settings differ across
designs, so these comparisons isolate neither lag nor report stability.
Longer delays remain to be evaluated.

\emph{Plots:} \path{multiqt_peak_reasoning_tradeoff.pdf}, pp.~1 and~7,
D (blended), three feedback blocks; \path{multiqt_evaluation_tree.pdf},
p.~1, COVID D in the delay-0, delay-1 and re-run branches. Flu D is
outside the ten methods on tree p.~7. For fixed rates, see
\path{multiqt_phase_rate_tradeoffs.pdf}, pp.~1, 7 and~13, $h=2$.

\subsection{Do quantile-specific expert weights help?}
For rate-grid mixtures, common weights lower pooled overall and
after-peak loss in every disease--feedback configuration, with or
without sharing and with an outer correction. Overall reductions are:
\begin{center}\small
\begin{tabular}{@{}lcc@{}}
\toprule
Rate-grid mixture & COVID: per-level $\to$ common & Flu: per-level $\to$ common\\
\midrule
No sharing & $-5.00\to-3.19$ & $-2.83\to-1.23$\\
Sharing & $-11.42\to-8.59$ & $-11.89\to-8.89$\\
Sharing + outer & $-11.74\to-8.91$ & $-12.17\to-9.16$\\
\bottomrule
\end{tabular}
\end{center}
The selected $H$ and outer rate match within pairs. Before the flu peak,
per-quantile weights without sharing give a small advantage at delay 0
and under re-run. Lower loss also need not mean lower CE.
\emph{Plots:} \path{multiqt_peak_reasoning_tradeoff.pdf}, pp.~1 and~7,
matched methods in all three windows; COVID forecast p.~2 and weight
p.~1 show US curves and weights at $h=2$, delay 0.

A--E use common weights only. Learned weekly starts have no matched
common-weight control, so their results do not establish a benefit from
per-quantile weighting. Equal weighting of all or recent starts changes
the weighting rule too. \emph{Plots:} COVID weight p.~2 (A--E) and
p.~3 (starts), with forecast/loss pp.~4 and~3.

\paragraph{Reading the plot references.}
Filenames and page numbers refer to the earlier \path{output/pdf/}
outputs, preserved under
\path{agent/outputs/notes/plots_history_20261111/} with a source manifest.
Forecast and weight files use the \path{multiqt_method_forecasts/} and
\path{multiqt_weight_timelines/} subdirectories. The all-location
trade-off markers are unweighted location means; tables and trees use
week-weighted pooling. Check horizon/design panels as well as pooled
results.

\paragraph{Loaded figures.}
\begin{center}\small
\renewcommand{\arraystretch}{1.12}
\begin{tabularx}{\textwidth}{@{}>{\raggedright\arraybackslash}Xll@{}}
\toprule
Original source PDF & Source pages & Notes start\\
\midrule
\path{multiqt_base_forecast_paths.pdf} & 1 & \pageref{plot:base-paths}\\
\path{multiqt_error_tradeoff.pdf} & 1, 4 & \pageref{plot:initial-tradeoff}\\
\path{multiqt_phase_rate_tradeoffs.pdf} & 1--7, 13 & \pageref{plot:phase-rates}\\
\path{multiqt_season_forecasts.pdf} & 1 & \pageref{plot:season-curves}\\
\path{multiqt_by_level.pdf} & 4, 12, 20, 27, 35, 43 & \pageref{plot:by-level}\\
\path{multiqt_tuning.pdf} & 1--6 & \pageref{plot:tuning}\\
\path{multiqt_peak_reasoning_tradeoff.pdf} & 1--12 & \pageref{plot:tradeoff}\\
\path{multiqt_evaluation_tree.pdf} & 1--12 & \pageref{plot:trees}\\
\path{multiqt_method_forecasts.pdf} & 2, 4 & \pageref{plot:method-curves}\\
COVID $h=2$, delay 0: studentized curves & 6, 7 & \pageref{plot:covid-studentized-curves}\\
Flu $h=2$, delay 0: studentized curves & 6, 7, 13, 14 & \pageref{plot:flu-studentized-curves}\\
\path{multiqt_weight_timelines.pdf} & 1, 2 & \pageref{plot:weights}\\
\path{multiqt_weight_timelines.pdf} & 3, 6, 9 & \pageref{plot:restart-weights}\\
\path{multiqt_method_forecasts.pdf} & 3, 10, 17 & \pageref{plot:restart-curves}\\
Flu $h=2$, delay 0: weights & 3, 6, 9 & \pageref{plot:flu-restart-weights}\\
Flu $h=2$, delay 0: forecasts & 3, 10, 17 & \pageref{plot:flu-restart-curves}\\
\path{multiqt_soft_switch_timelines.pdf} & 1 & \pageref{plot:switch}\\
\bottomrule
\end{tabularx}
\end{center}
The singleton forecast, weight and switch PDFs are COVID, $h=2$,
delay 0; flu uses \path{flu_horizon2_delay0.pdf} in the matching
subdirectory. Tuning pages show COVID delay 0, delay 1 and re-run,
then the same flu designs, under the earlier procedure. The season
and by-level pages are exact extracts from \path{notes_before_update.pdf}
in the archive. These frozen figures are not replaced by new workflow
outputs; updating the results requires replacing scores, figures and
page references together. Notes-only copies change wording, not data.

\begin{thebibliography}{7}
\bibitem{multiqt}
Tiffany Ding, Isaac Gibbs, and Ryan J. Tibshirani.
\emph{Calibrated Multi-Level Quantile Forecasting}.
arXiv:2512.23671v1, 2025.
\href{https://arxiv.org/abs/2512.23671v1}{MultiQT paper}.

\bibitem{pid}
Anastasios N. Angelopoulos, Emmanuel J. Cand\`es, and Ryan J. Tibshirani.
\emph{Conformal PID Control for Time Series Prediction}.
Advances in Neural Information Processing Systems 36, 2023.
\href{https://proceedings.neurips.cc/paper_files/paper/2023/file/47f2fad8c1111d07f83c91be7870f8db-Paper-Conference.pdf}{NeurIPS paper}.

\bibitem{delphi-calibration}
CMU Delphi.
\emph{Delphi hub forecasting: calibration experiments and tracker code}.
Pinned source revision \texttt{9abebcfb}, inspected 9--10 November 2026.
\href{https://github.com/cmu-delphi/delphi-hub-forecasting/blob/9abebcfb/reports/writeups/calibration_experiments/e09_lr_delay.Rmd}{E09: constant learning rate and revision delay};
\href{https://github.com/cmu-delphi/delphi-hub-forecasting/blob/9abebcfb/reports/writeups/calibration_experiments/e15_wis_sources.Rmd}{E15: median shift and width change};
\href{https://github.com/cmu-delphi/delphi-hub-forecasting/blob/9abebcfb/R/calibration/qt.R}{tracker implementation}.

\bibitem{gneiting2005}
Tilmann Gneiting, Adrian E. Raftery, Anton H. Westveld III, and Tom Goldman.
\emph{Calibrated Probabilistic Forecasting Using Ensemble Model Output
Statistics and Minimum CRPS Estimation}.
Monthly Weather Review 133(5):1098--1118, 2005.
\href{https://doi.org/10.1175/MWR2904.1}{Paper}.

\bibitem{gibbs-candes2024}
Isaac Gibbs and Emmanuel J. Cand\`es.
\emph{Conformal Inference for Online Prediction with Arbitrary Distribution
Shifts}.
Journal of Machine Learning Research 25(162):1--36, 2024.
\href{https://jmlr.org/papers/v25/22-1218.html}{Paper and implementation}.

\bibitem{berrisch-ziel2023}
Jonathan Berrisch and Florian Ziel.
\emph{CRPS Learning}.
Journal of Econometrics 237, article 105221, 2023.
\href{https://doi.org/10.1016/j.jeconom.2021.11.008}{Journal article};
\href{https://arxiv.org/abs/2102.00968}{Author manuscript}.

\bibitem{expert-design-records}
Project experiment records, October 2026.
Earlier rate-bank diagnosis: \path{agent/outputs/postpeak_diagnosis/REPORT.md}.
Experiments and plan of 28 October:
\path{agent/outputs/peak_antiwindup/README.md} and
\path{agent/docs/peak_correction_todos.md}.
Diverse-bank experiment of 30 October:
\path{agent/outputs/expert_bank_comparison/README.md}, with saved
forecasts and weights in \path{forecasts.rds} in the same directory.
These are historical exploratory records, not the current evaluation protocol.
\end{thebibliography}

\clearpage
\appendix
\section{Loaded plot pages}
\label{sec:loaded-plots}


\includepdf[pages=1,fitpaper=true,pagecommand={\thispagestyle{empty}},
 pagecommand*={\thispagestyle{empty}\phantomsection\label{plot:base-paths}}]{../agent/outputs/notes/plots_history_20261111/multiqt_base_forecast_paths.pdf}
\includepdf[pages={1,4},fitpaper=true,pagecommand={\thispagestyle{empty}},
 pagecommand*={\thispagestyle{empty}\phantomsection\label{plot:initial-tradeoff}}]{../agent/outputs/notes/plots_history_20261111/multiqt_error_tradeoff.pdf}
\includepdf[pages={1-7,13},fitpaper=true,pagecommand={\thispagestyle{empty}},
 pagecommand*={\thispagestyle{empty}\phantomsection\label{plot:phase-rates}}]{../agent/outputs/notes/plots_history_20261111/multiqt_phase_rate_tradeoffs.pdf}
\includepdf[pages=1,fitpaper=true,pagecommand={\thispagestyle{empty}},
 pagecommand*={\thispagestyle{empty}\phantomsection\label{plot:season-curves}}]{../agent/outputs/notes/plots_history_20261111/multiqt_season_forecasts_selected.pdf}
\includepdf[pages={1-6},fitpaper=true,pagecommand={\thispagestyle{empty}},
 pagecommand*={\thispagestyle{empty}\phantomsection\label{plot:by-level}}]{../agent/outputs/notes/plots_history_20261111/multiqt_by_level_selected.pdf}
\includepdf[pages={1-6},fitpaper=true,pagecommand={\thispagestyle{empty}},
 pagecommand*={\thispagestyle{empty}\phantomsection\label{plot:tuning}}]{../agent/outputs/notes/plots_history_20261111/multiqt_tuning.pdf}
\includepdf[pages={1-12},fitpaper=true,pagecommand={\thispagestyle{empty}},
 pagecommand*={\thispagestyle{empty}\phantomsection\label{plot:tradeoff}}]{../agent/outputs/notes/plots_history_20261111/multiqt_peak_reasoning_tradeoff.pdf}
\includepdf[pages={1-12},fitpaper=true,pagecommand={\thispagestyle{empty}},
 pagecommand*={\thispagestyle{empty}\phantomsection\label{plot:trees}}]{../agent/outputs/notes/plots_history_20261111/multiqt_evaluation_tree.pdf}
\includepdf[pages={2,4},fitpaper=true,pagecommand={\thispagestyle{empty}},
 pagecommand*={\thispagestyle{empty}\phantomsection\label{plot:method-curves}}]{../agent/outputs/notes/plots_history_20261111/multiqt_method_forecasts.pdf}
\includepdf[pages={6,7},fitpaper=true,pagecommand={\thispagestyle{empty}},
 pagecommand*={\thispagestyle{empty}\phantomsection\label{plot:covid-studentized-curves}}]{../agent/outputs/notes/plots_history_20261111/multiqt_method_forecasts.pdf}
\includepdf[pages={6,7,13,14},fitpaper=true,pagecommand={\thispagestyle{empty}},
 pagecommand*={\thispagestyle{empty}\phantomsection\label{plot:flu-studentized-curves}}]{../agent/outputs/notes/plots_history_20261111/multiqt_method_forecasts/flu_horizon2_delay0.pdf}
\includepdf[pages={1,2},fitpaper=true,pagecommand={\thispagestyle{empty}},
 pagecommand*={\thispagestyle{empty}\phantomsection\label{plot:weights}}]{../agent/outputs/notes/plots_history_20261111/multiqt_weight_timelines.pdf}
\includepdf[pages={3,6,9},fitpaper=true,pagecommand={\thispagestyle{empty}},
 pagecommand*={\thispagestyle{empty}\phantomsection\label{plot:restart-weights}}]{../agent/outputs/notes/plots_history_20261111/multiqt_weight_timelines.pdf}
\includepdf[pages={3,10,17},fitpaper=true,pagecommand={\thispagestyle{empty}},
 pagecommand*={\thispagestyle{empty}\phantomsection\label{plot:restart-curves}}]{../agent/outputs/notes/plots_history_20261111/multiqt_method_forecasts.pdf}
\includepdf[pages={3,6,9},fitpaper=true,pagecommand={\thispagestyle{empty}},
 pagecommand*={\thispagestyle{empty}\phantomsection\label{plot:flu-restart-weights}}]{../agent/outputs/notes/plots_history_20261111/multiqt_weight_timelines/flu_horizon2_delay0.pdf}
\includepdf[pages={3,10,17},fitpaper=true,pagecommand={\thispagestyle{empty}},
 pagecommand*={\thispagestyle{empty}\phantomsection\label{plot:flu-restart-curves}}]{../agent/outputs/notes/plots_history_20261111/multiqt_method_forecasts/flu_horizon2_delay0.pdf}
\includepdf[pages=1,fitpaper=true,pagecommand={\thispagestyle{empty}},
 pagecommand*={\thispagestyle{empty}\phantomsection\label{plot:switch}}]{../agent/outputs/notes/plots_history_20261111/multiqt_soft_switch_timelines.pdf}

\end{document}
